% !TeX root = ../../Book3.tex
% 纲要标题：### 19.4 超曲面、嵌入分支与重数
% 纲要位置：计划纲要.md，第 2106 行。

\section{超曲面、嵌入分支与重数}
\label{b3:sec:1904}

一个超曲面可以约化或非约化、正规或不正规，却仍然是 Cohen–Macaulay 环。将这些例子与具有嵌入关联素理想的商环比较，可以区分零因子、奇点与深度不足各自造成的现象。


% 纲要内容：
%
% * 超曲面与零维局部环
% * 非 Cohen–Macaulay 的嵌入分支例子
% * 参数、重数与几何交的联系
%
% ---
%

\subsection{超曲面与零维局部环}
\label{b3:subsec:1904:01}

设 \((S,\mathfrak n)\) 为交换 Noether 正则局部环，\(0\ne f\in\mathfrak n\)。商 \(S/(f)\) 称为这里的\term[chaoqumianhuan]{超曲面环}[hypersurface ring]。因为 \(S\) 为整环，\(f\) 是非零因子，所以超曲面环是完全交环，深度和维数都等于 \(\dim S-1\)。是否约化、是否正规、是否正则，则取决于方程的进一步性质。

\ProseParagraphBreak
在 \(S=k[[x,y]]\) 中，\(S/(y^2-x^3)\)、\(S/(xy)\) 和 \(S/(x^2)\) 都是一维 CM 环。第一个在代数上有单个尖点分支，第二个有两个相交分支，第三个含非零幂零元。CM 性要求存在参数非零因子，没有把这些不同现象合并为一种光滑性。

\begin{proposition}\label{b3:prop:zero-dimensional-cm}
每个零维交换 Noether 局部环及其非零有限模都是 CM 的。
\end{proposition}
\begin{proof}
设环的极大理想为 \(\mathfrak m\)。它是幂零理想，所以非零有限模 \(M\) 有有限长度。取最小的 \(r\ge1\) 使 \(\mathfrak m^rM=0\)，则 \(\mathfrak m^{r-1}M\ne0\) 被 \(\mathfrak m\) 零化。因而不存在属于 \(\mathfrak m\) 的 \(M\) 非零因子，深度为零；模维数也为零。
\end{proof}

零维时 CM 性没有额外限制，完全交性却仍有限制；上一节的基座计算就是一种有效区分。

\subsection{非 Cohen–Macaulay 的嵌入分支}
\label{b3:subsec:1904:02}

令
\[
A=k[x,y]_{(x,y)}/(x^2,xy),\qquad\mathfrak m=(\bar x,\bar y).
\]
其根理想为 \((x)\)，故 \(\dim A=1\)。又 \(\bar x\ne0\) 且
\(\mathfrak m\bar x=0\)，于是 \(\mathfrak m\in\operatorname{Ass}_AA\) 且 \(\operatorname{depth}A=0\)。这里既有一条一维分支，也有集中在原点的零化关系。它不是只有“图形上看起来多了一个点”这么简单：乘法 \(\bar y\) 的核恰好暴露了这个额外部分。

\ProseParagraphBreak
\(\bar y\) 是参数，因为 \(A/(\bar y)\cong k[x]/(x^2)\) 长度为二；但它不是非零因子。若用一般参数 \(\bar y+a\bar x\)，仍有 \((\bar y+a\bar x)\bar x=0\)。所以换一个参数并不能修复深度不足，参数系判据要求的正则性确实失败。

\ProseParagraphBreak
一维约化 Noether 局部环则为 CM 环。其零因子是极小素理想的并：由零理想为这些素理想的交可见，避开所有极小素理想的乘法是单射；反向可由关联素理想判据得到。有限素理想回避使极大理想中存在这样的元素，深度至少为一；维数上界再给出深度等于一。到了二维，“约化”只能保证第一个非零因子，不能保证第二个。

\subsection{参数、重数与几何交}
\label{b3:subsec:1904:03}

参数理想给出的有限长度保留了相交时的重数信息。一个最直接、又不需要一般相交理论的计算如下。

\begin{proposition}\label{b3:prop:cm-parameter-powers-length}
设 \((R,\mathfrak m)\) 为交换 Noether 局部环，\(M\ne0\) 为 \(d\) 维有限 CM 模，\(x_1,\ldots,x_d\) 为参数系。对正整数 \(a_1,\ldots,a_d\)，有
\[
\ell_R\bigl(M/(x_1^{a_1},\ldots,x_d^{a_d})M\bigr)
=a_1\cdots a_d\,\ell_R\bigl(M/(x_1,\ldots,x_d)M\bigr).
\]
\end{proposition}
\begin{proof}
幂次参数系仍有同样的根理想，故仍为参数系；\cref{b3:thm:cm-parameters}使它及其任意重排均正则。先把除 \(x_1\) 以外的参数幂商去，得模 \(N\)，其中 \(x_1\) 为非零因子。对每个 \(j\ge0\)，乘 \(x_1^j\) 诱导同构
\[
N/x_1N\longrightarrow x_1^jN/x_1^{j+1}N.
\]
满射来自定义；若 \(x_1^jn=x_1^{j+1}n'\)，可约去非零因子 \(x_1^j\)，故核为零。用这些商滤过 \(N/x_1^{a_1}N\)，长度相加得到因子 \(a_1\)。依次对其余参数做同样的计算。
\end{proof}

在 \(k[[x,y]]\) 中，参数理想 \((x^a,y^b)\) 的商以
\(x^iy^j\ (0\le i<a,\ 0\le j<b)\) 为基，长度为 \(ab\)。它的支集只有原点，但商代数记住了两条带重数坐标轴相交的程度。

\ProseParagraphBreak
一维时还能直接把长度与 Hilbert–Samuel 重数比较。若 \(M\) 为一维 CM 模，\(t\) 为参数，则上述滤过给出
\(\ell(M/t^{n+1}M)=(n+1)\ell(M/tM)\)，故
\(e_{(t)}(M)=\ell(M/tM)\)。例如尖点环
\(R=k[[x,y]]/(y^2-x^3)\) 中，\(x\) 为非零因子，\(R/(x)\cong k[y]/(y^2)\)，因而 \(e_{(x)}(R)=2\)。

\ProseParagraphBreak
上一小节的非 CM 环 \(A\) 则有
\(\ell(A/(\bar y^{n+1}))=n+2\)：可取基
\(1,\bar y,\ldots,\bar y^n,\bar x\)。因此 \(e_{(\bar y)}(A)=1\)，而
\(\ell(A/(\bar y))=2\)。被 \(\bar y\) 零化的那一维部分增加常数项，却不增加一维增长的首项。这个差别说明参数商长度与重数的相等需要实质条件。
